\section{End-to-end cost of one fixed-source phase}
\label{sec:complexity}

\subsection{Compilation, observations and monotone updates}

Write $P(k,m,B)$ for a valid polynomial bound on source-profile compilation
through weighted AHT, including decoding and exact source validation.
Let $s$ be the number of profile rows, $e$ its support incidence count,
and let $Q$ later observations each have dimension at most $d$ and weight
entries of at most $D$ bits. A conservative batch bound is
\begin{equation}
 P(k,m,B)+O\bigl(Qsmd\,\mathsf M(B+D+\log(m+1))\bigr)
 \label{eq:batch-cost}
\end{equation}
plus canonical histogram grouping. This replaces $Q$ separate weighted
trace transports by one compilation and $Q$ explicit linear maps on the
table. It can lose for a small number of low-dimensional queries, because
compiling all source coordinates may cost more than answering them directly.
No unconditional speed ordering between these strategies is asserted.

For an explicit monotone cone sequence, define the index work
\[
 U(s,m,e,u)=(s+m+e+u)\alpha(s+m)+u\log(m+1).
\]
The online count bound is $O(smB)$ for reading the already validated dense
coordinates, followed by $O(U)$ index operations and at most
$O(s+m)$ effective count changes with $B$-bit arithmetic. A simple all-bit
upper bound charges $O(\log(s+m+1)+B)$ to every index-level operation;
it is intentionally conservative. Census output, when requested, adds
$O(sm\,\mathsf M(B))$ arithmetic and sorting. Repeatedly requesting a full
dense census is output work, not a violation of the online count bound.

The row bound in \cref{thm:quotient} is useful here. Within a fixed-source
phase, quotient profiles never require more than the original $s$ rows,
and every coordinate remains at most its original atom capacity. Thus
neither histogram length nor coordinate bit size grows through these
quotients. This statement concerns the chosen observer, not the size of
future geometric descriptions.

\subsection{Compressed schedule accounting}

Let $T_{\rm trial}$ bound cloning a quotient state and applying one stored
summary. A safe choice in index operations is
\[
 T_{\rm trial}=O\bigl((s+m+e)\alpha(s+m)+m\log(m+1)\bigr),
\]
with the preceding count-bit charges. A summary has at most $m$ atom
entries, so it never contains an expanded program. The immutable adjacency
lists are shared across clones, and their initial $O(sm)$ construction is
not repeated by a clone.

After source compilation and grammar preparation, the first-threshold
algorithm costs $O(gT_{\rm trial}+gH)$ in the separated model, while all
history events cost
\[
 O\bigl((h+1)g(T_{\rm trial}+H)\bigr),\qquad h\le2m-1.
\]
Full grammar preparation additionally pays for its explicit leaf occurrences,
canonical summary joins, and $g$ binary length computations. There is no
factor equal to the numerical expanded length $L(P)$.

The independent checker has its own polynomial cost. Its graph-based summary
join can use quadratic atom work, and it may rebuild the grammar data for
several prefix queries. These are deliberate simplicity choices, not costs
silently charged to the producer. A host measuring a certified operation
must include both generation and replay. The timing tables below measure
production search, with answer checks outside the timed region, and do not
label those times as certified end-to-end native performance.

\begin{corollary}[Polynomial fixed-source observer]
\label{cor:local-poly}
A binary interval-pairing source, a fixed atom partition, and a finite
straight-line program of full-port cones admit exact final census,
first-count-threshold and complete count-change history computations in
time polynomial in their encoded input sizes. Corresponding algebra-level
claims admit polynomial independent replay. A source-bound replay additionally
requires a verified weighted orbit certificate.
\end{corollary}

\begin{proof}
Compilation is polynomial by weighted AHT and \cref{thm:packing}.
Its compressed output has polynomial $s$ and hence polynomial $e\le sm$.
All arithmetic uses polynomially many bits, and the grammar's length bits
are bounded by its encoded exponents and node count. Apply the preceding
bounds and \cref{thm:history,prop:history-cert}.
\end{proof}

\subsection{How this enters a quasi-polynomial programme}

Consider a prospective recognizer divided into $R$ fixed-source phases.
For phase $i$, record its source parameters, actual compilation cost bound,
number of observation coordinates, grammar size, and event count. Let
$\mathcal G$ denote all work outside this observer: source generation,
normal-vector discovery, actual topological cuts, verification of attachment
semantics, and transfers between phases. The honest total ledger is
\begin{align}
 T_{\rm total}\ \le\ \mathcal G+
 \sum_{i=1}^{R}\bigl[&P(k_i,m_i,B_i)+\cost(\text{grammar}_i)
 +\cost(\text{observations}_i)\notag\\
 &+O((h_i+1)g_i(T_{{\rm trial},i}+H_i))
 +\cost(\text{replay}_i)\bigr].
 \label{eq:global-ledger}
\end{align}
If $R$ and all encoded phase descriptions are at most
$2^{O((\log n)^c)}$, for fixed $c$, then the displayed observer sum is
quasi-polynomial in the crossing input size $n$. This is a consequence of
polynomial closure of that size bound, not a proof that the required
phases exist. In particular, \eqref{eq:global-ledger} says nothing favorable
about $\mathcal G$ without additional theorems.

There is nevertheless a concrete strengthening over per-instruction
accounting. A phase with exponentially many redundant fixed-port operations
can be processed through its grammar and at most $2m-1$ effective changes.
This removes its expanded instruction count from the ledger. A mere
termination proof based on a binary counter could not make that replacement.
The achievement is local and structural: a small-height quotient algebra
has replaced a potentially huge operation sequence.

\subsection{Why no global recognition theorem follows}

Several missing links are mathematically separate. A supplied normal surface
can be analyzed without finding a useful one. A valid triangulation can be
analyzed without proving that it is the exterior of the input knot. A
full-port coning operation can be evaluated without proving that it models
a particular hierarchy move. Finally, a bounded-height phase can terminate
quickly while the number of phases or the required source refinements is
large.

The four-point example proves that one cannot simply place all future
attachment maps behind this interface. The fixed-atom potential also resets
when fresh atoms or changed source relations are introduced. A useful
continuation should therefore target a specific geometric phase in which
these semantics can actually be proved, and then measure how often that
phase is rebuilt. Treating a local polynomial observer as a complete
quasi-polynomial recognizer would skip precisely the hard accounting that
\eqref{eq:global-ledger} makes explicit.
